% =============================================================================
% slides_rebonato_tips_v3.tex — Meeting deck, Paper 2 (TIPS–Treasury Floor)
% PREAMBLE: copiato VERBATIM da slides/presentation.tex (CambridgeUS, 16:9,
% divisori di sezione rossi via \AtBeginSection). Contenuto: versione light.
% =============================================================================
\documentclass[aspectratio=169]{beamer}

\mode<presentation> {
	\usetheme{CambridgeUS}
}

\AtBeginSection[]{
  \begin{frame}
  \vfill
  \centering
  \begin{beamercolorbox}[sep=8pt,center,shadow=true,rounded=true]{title}
    \usebeamerfont{title}\textcolor{red}\insertsectionhead\par%
  \end{beamercolorbox}
  \vfill
  \end{frame}
}

\IfFileExists{transparent.sty}{\usepackage{transparent}}{}
\IfFileExists{dsfont.sty}{\usepackage{dsfont}}{}
% \usepackage{stackengine}   % riattivare in locale (qui manca la dipendenza listofitems)
% \newcommand\barbelow[1]{\stackunder[1.4pt]{$#1$}{\rule{.8ex}{.075ex}}}   % non usato in questo deck
\setbeamertemplate{navigation symbols}{}
\usepackage{graphicx}
\usepackage{booktabs}
\IfFileExists{kotex.sty}{\usepackage{kotex}}{}
\usepackage{amsmath, amsthm, amssymb, caption, bm}
\newtheorem{prop}{Proposition}

\newcommand{\bps}{\,bp}
\newcommand{\VaR}{\mathrm{VaR}}
\newcommand{\CE}{\mathrm{CE}}

%----------------------------------------------------------------------------------------
%	TITLE PAGE
%----------------------------------------------------------------------------------------

\title[]{The TIPS--Treasury Basis as a\\ Limits-to-Arbitrage Equilibrium Floor}
\author{Alessio Ottaviani}
\institute[EDHEC Business School]{EDHEC Business School\\ \smallskip Supervisor: Prof.\ Riccardo Rebonato}
\date{July 2026}

\begin{document}

%%%%%%%%%%%%%%%
\begin{frame}
    \titlepage
\end{frame}
%%%%%%%%%%%%%%%

\begin{frame}{Agenda}
    \tableofcontents
\end{frame}

%----------------------------------------------------------------------------------------
\section{The Idea}
%----------------------------------------------------------------------------------------
\begin{frame}{\emph{The Idea}, restated}
\begin{itemize}\setlength\itemsep{0.5em}
  \item TIPS $+$ zero-coupon inflation swap replicates a nominal Treasury; the synthetic
        trades \emph{cheap} by $\lambda_t$: a textbook LOOP violation.
  \item The convergence trade has \textbf{certain} terminal payoff
        $S(T)=1-e^{-\lambda_0 T}>0$ --- \emph{if it can be held}.
  \item Before $T$, mark-to-market losses can breach the risk limit:
        $\tau_B=\inf\{t:\ D(\lambda_t-\lambda_0)\ge B\}$, $B=m\times\VaR_\alpha$
        --- forced unwinding at the worst moment.
  \item Entry iff $\CE\!\left[W_0+X_T\right]\ge W_0$, with $W_0=$ capital charge;
        free entry pins the \textbf{equilibrium floor} $\lambda^*$.
\end{itemize}
\vspace{0.3em}
\begin{block}{Claim}
The post-compression basis is not a dying anomaly --- it is $\lambda^*$, observed.
This deck: the question, first evidence, and the plan.
\end{block}
\end{frame}

%----------------------------------------------------------------------------------------
\section{First Evidence}
%----------------------------------------------------------------------------------------
\begin{frame}{Twenty years of the mispricing, on FLL's own construction}
\begin{center}
\includegraphics[width=0.82\textwidth]{fll_fig2_extended.png}
\end{center}
\vspace{-0.4em}
{\small Bond-level, 107 CUSIP pairs, notional-weighted (their Figure 2, replicated).
\textbf{Everything to the right of the dashed line is the out-of-sample of
Fleckenstein--Longstaff--Lustig (2014, JF)} --- and the object of this paper.}
\end{frame}

\begin{frame}{The floor exists: 16--24\bps, fifteen years after FLL}
\begin{center}\small
\begin{tabular}{lcccc}
\toprule
Monthly mean [NW $t$] & 2Y & 5Y & 10Y & 30Y \\
\midrule
FLL window 04/07--09/11 & 44.0\,[2.9] & 34.6\,[9.2] & \textbf{46.8\,[7.5]} & 51.0\,[8.5] \\
Post-2013 & 23.3\,[7.6] & 15.7\,[10.9] & \textbf{23.4\,[22.4]} & 24.0\,[13.3] \\
\bottomrule
\end{tabular}
\end{center}
\vspace{0.4em}
\begin{itemize}
  \item Construction validated three ways: breakeven cross-check (corr.\ 0.996 at 10Y),
        bond-level median (0.926), independent parallel reconstruction.
  \item Post-2013: \emph{stationary} (ADF $p=0.008$), OU half-life 3.7--4.6 months around
        $\theta\approx17\bps$ --- mean reversion \emph{around a positive level}.
  \item End effects quantified: 2Y carries a 52\bps\ CPI seasonal; 20Y (from 2020) is the
        sector-reintroduction cheap effect.
\end{itemize}
\end{frame}

\begin{frame}{Compression is a staircase, not a decay}
\begin{center}
\includegraphics[width=0.66\textwidth]{fig7_breaks10y.png}
\end{center}
\vspace{-0.3em}
{\small Ten-year: dominant break \textbf{December 2010} (sup-$F=160$), second break
\textbf{December 2017} ($46\to29\to21\bps$). Since October 2024 the 2Y basis is
statistically zero: the short-end arbitrage has closed. Next step in the plan:
line these dates up against the regulatory capital calendar (H4).}
\end{frame}

\begin{frame}{The barrier signature: violent conditionally, empty linearly}
\begin{columns}[T]
\begin{column}{0.48\textwidth}
\textbf{Unconditional linear: nothing}
\begin{itemize}
  \item Local projections on MOVE shocks: $|t|\le1.2$ at every horizon.
  \item Rolling resting level on rolling MOVE: $\beta=-0.014$ $[t=-0.4]$.
\end{itemize}
\end{column}
\begin{column}{0.48\textwidth}
\textbf{Conditional: 3--7$\times$ moves, with lags}
\begin{itemize}
  \item GFC 144--154\bps; Mar-2020 (2Y) $-1\to68$; 2022 $\to68$; SVB $\to60$.
  \item Peaks \emph{lag} the shock by weeks--months (SVB: March shock, July--August peaks).
\end{itemize}
\end{column}
\end{columns}
\vspace{0.7em}
\begin{alertblock}{Reading}
A linear regression averages a flat normal-times relation with a violent
constraint-binding one --- and recovers nothing. The signature of a constraint that
binds occasionally.
\end{alertblock}
\end{frame}

\begin{frame}{A first, model-consistent sorting of the episodes}
\begin{center}\small
\begin{tabular}{lccc}
\toprule
Episode onset & HKM capital ratio & Percentile (2004+) & Bottom quintile \\
\midrule
GFC (Oct-2008) & 3.57\% & \textbf{3\%} & yes \\
COVID (Mar-2020) & 4.27\% & \textbf{8\%} & yes \\
2022 (Jun-2022) & 5.41\% & 32\% & no \\
SVB (Mar-2023) & 5.62\% & 37\% & no \\
\bottomrule
\end{tabular}
\end{center}
\vspace{0.5em}
\begin{itemize}
  \item The two \emph{capital-impairment} onsets are exactly the two system-wide blow-outs;
        the unimpaired onsets produced the smaller, short-end moves.
  \item In March 2020 the bond-level \emph{cross-section fans out before the median moves}
        (IQR peaks on the dash-for-cash apex, 16 March) --- the heterogeneous-threshold
        prediction.
\end{itemize}
\vspace{0.2em}
{\small These patterns motivate the formal program on the next slide.}
\end{frame}

%----------------------------------------------------------------------------------------
\section{Research Plan}
%----------------------------------------------------------------------------------------
\begin{frame}{Research plan}
\begin{enumerate}\setlength\itemsep{0.3em}
  \item \textbf{State-dependent constraint tests}: threshold responses conditioned on the
        \emph{lagged constraint state} (HKM; FR\,2004 dealer positions; CFTC), lag structure
        from the event studies, bootstrap inference. Known risk, stated upfront: two U.S.
        impairment clusters $\Rightarrow$ power --- hence (3).
  \item \textbf{The structural floor}: estimate the dynamics, compute the first-passage
        survival law, solve the CE entry condition for $\lambda^*$ --- with the capital
        charge \emph{derived} from $m\times\VaR$ under the estimated dynamics
        (regulatory counterfactuals on $(m,\alpha)$ for free).
  \item \textbf{The international cross-section}: add markets supplying episodes the U.S.\
        sample lacks. \emph{Euro} (Bund\texteuro i-vs-OAT\texteuro i differential, quanto-CDS
        controls): the 2011--12 crisis, sovereign risk held fixed. \emph{UK} index-linked
        gilts: the Oct-2022 LDI dislocation --- orthogonal to both U.S.\ clusters, no
        sovereign confound, no deflation floor.
  \item \textbf{Structural estimation (SMM)}: survival surface by vintage $\times$ buffer
        $+$ state-dependent loadings as over-identifying moments; block-bootstrap moment
        covariance.
\end{enumerate}
\end{frame}

%----------------------------------------------------------------------------------------
\section{Discussion}
%----------------------------------------------------------------------------------------
\begin{frame}{Three points where I need your judgment}
\begin{enumerate}\setlength\itemsep{0.8em}
  \item \textbf{The international extension}: is the Bund\texteuro i-vs-OAT\texteuro i
        differential, with quanto-CDS controls, the identification you would trust for
        2011--12 --- and would you place the UK gilt basis (Oct-2022 LDI) ahead of it,
        or alongside?
  \item \textbf{Modelling the capital charge} (Phase 2): standalone versus \emph{marginal}
        VaR --- the question your note closes on --- and the normalisation of risk aversion
        in the CE condition (curvature over committed capital vs total wealth).
  \item \textbf{Ambition}: top field tier on the strength of the first evidence, with the
        general-interest route conditional on Phases 2--3 --- and how you would like us to
        work on this together.
\end{enumerate}
\end{frame}

%----------------------------------------------------------------------------------------
\appendix
\section*{Backup}
%----------------------------------------------------------------------------------------
\begin{frame}{Backup --- Construction validation}
\begin{center}\small
\begin{tabular}{lcccc}
\toprule
Maturity & Mean diff vs BE (bp) & SD & Corr. & $n$ \\
\midrule
2Y & $-3.97$ & 9.78 & 0.974 & 5648 \\
5Y & $-5.88$ & 9.76 & 0.839 & 5718 \\
10Y & $-0.55$ & 1.55 & \textbf{0.996} & 5718 \\
30Y & $+0.60$ & 4.70 & 0.969 & 5718 \\
\bottomrule
\end{tabular}
\end{center}
{\small Direct construction vs Bloomberg-breakeven route, daily. Bond-level median:
correlation 0.926, common regimes, common peaks; the cross-sectional average replicates
FLL's Figure 2 on their window.}
\end{frame}

\begin{frame}{Backup --- March 2020 at bond level}
\begin{center}
\includegraphics[width=0.62\textwidth]{fig10_mar2020_bondlevel.png}
\end{center}
{\small The median doubles (16 March) while the IQR triples, peaking on the dash-for-cash
apex; the median's own maximum lags to late June. The cross-section moves first ---
heterogeneous thresholds at work.}
\end{frame}

\end{document}
